\documentclass[a4paper,10pt]{article}
\usepackage[margin=1.5cm]{geometry}
\usepackage{multicol}
\usepackage{hyperref}
\usepackage{enumitem}
\usepackage{listings}



\setlength{\parindent}{0pt}

\begin{document}



\begin{center}
    {\Large \texttt{Linux Trouble Shooting}}\\
\end{center}

\tableofcontents
\newpage



\section{Introducción}

En este documento encontrarás los soluciones a los problemas del kernel/os de distribuciones \textit{Linux} con los que me he encontrado.



\section{Kernel Panic}


\subsection{VFS: Unable to mount root fs on unkowun-block 0,0}

Este error ocurre cuando hubo un problema al momento de compilar el kernel y el \textit{initramfs}. Este último es una imagen temporal del filesystem que se usa para ayudar al kernel a montar el root filesystem y correr el sistema principal \textit{Init}.
En mi caso, el error se debió a que un modulo de kernel de terceros, \textit{hid-xpadneo},  no era compatible con la versión de kernel con la que traté de iniciar mi equipo, Linux 7.0.0-15-generic.
Para encontrar qué módulo fue el que falló, corrí el siguiente comando:

\begin{lstlisting}
    sudo dkms autoinstall -k 7.0.0-15-generic (cambia por la version de kernel del caso)
\end{lstlisting}

En los logs de ejecución, se podía ver que había ocurrido un error al tratar de integar el paquete hid-xpadneo en la versión de kernel. En mi caso, ese modulo ya no era necesario, por lo que la solución fue \textbf{removerlo}.
Para saber correctamente el nombre del paquete, se puede consultar el siguiente comando:



\begin{lstlisting}
    sudo dkms status
\end{lstlisting}

El output de ese comando mostará los paquetes de terceros presentes para ser integrados al kernel. Los estatus pueden ser "added" o "installed". Incluso si se muestra solo como agregado, al momento de tratar de compilar al kernel también mostrará el error. El objetivo es que el paquete desaparezca. Para ello, ejecuta:

\begin{lstlisting}
    sudo dkms remove hid-xpadneo/v0.9-214-g8d20a23(cambia por el paquete problematico) --all 
\end{lstlisting}

Una vez ejecutado eso, corroboramos con dkms status. Si se tuvo exito, el paquete ya no aparecerá en el output de ese comando.


Ahora, para \textbf{reparar} la compilación del kernel, sigue la secuencia de comando:

\begin{lstlisting}
    sudo apt update
    sudo apt upgrade
    sudo dkms autoinstall -k 7.0.0-15-generic  
    sudo update-initramfs -u -k 7.0.0-15-generic
\end{lstlisting}

Después de eso, reinicia el equipo e intenta iniciar con el kernel reparado.


\end{document}